% Aqui começa o capítulo de Tailoring do Projeto.
\chapter{Tailoring do Projeto}\label{chp:Tailoring}

O PMBOK 7ª edição enfatiza que não existe uma abordagem única e universal para o gerenciamento de projetos. O conceito de \textit{tailoring} refere-se à adaptação consciente de princípios, práticas, métodos e ferramentas às características específicas do projeto, da organização e do ambiente em que ele está inserido. Neste capítulo, são apresentadas as decisões de tailoring adotadas no projeto de desenvolvimento do sistema de gestão de casos sociais, justificando a escolha e a adaptação do framework Scrum.

\section{Justificativa da Abordagem Adotada}

O projeto de gestão de casos sociais apresenta características que tornam inadequada a adoção de uma abordagem puramente preditiva. Entre essas características destacam-se a incerteza dos requisitos, a diversidade de stakeholders, a influência de fatores humanos e organizacionais e a necessidade de validação contínua da solução com usuários reais.

Diante desse contexto, optou-se pela adoção do Scrum como framework ágil principal, por sua capacidade de lidar com complexidade, promover entregas incrementais de valor e favorecer o aprendizado contínuo. Entretanto, a aplicação do Scrum não ocorre de forma estrita, sendo ajustada conforme as restrições institucionais e os objetivos do projeto, em consonância com o conceito de tailoring proposto pelo PMBOK 7.

\section{Adaptação dos Papéis do Scrum}

Os papéis definidos pelo Scrum foram mantidos, porém adaptados à realidade do projeto simulado. O Product Owner é responsável por representar os interesses institucionais e dos usuários finais, atuando como principal elo entre os stakeholders e a equipe de desenvolvimento. Esse papel incorpora também a responsabilidade de equilibrar demandas sociais, legais e operacionais na priorização do backlog.

O Scrum Master atua como facilitador do processo, promovendo a adoção das práticas ágeis, a remoção de impedimentos e a melhoria contínua da forma de trabalho. Considerando o contexto acadêmico e institucional do projeto, esse papel também assume a função de promover alinhamento metodológico entre os integrantes do grupo.

A equipe de desenvolvimento é formada de maneira multifuncional, reunindo competências técnicas e analíticas necessárias para a concepção do sistema. A auto-organização da equipe é incentivada, respeitando as limitações de tempo e recursos inerentes ao projeto.

\section{Adaptação dos Eventos do Scrum}

Os eventos do Scrum foram ajustados para garantir equilíbrio entre disciplina e flexibilidade. As sprints foram definidas com duração fixa, permitindo ciclos curtos de entrega e aprendizado, sem comprometer a previsibilidade do planejamento.

O planejamento da sprint é utilizado para definir objetivos claros e selecionar itens do backlog compatíveis com a capacidade da equipe. As reuniões diárias são realizadas de forma objetiva, com foco na sincronização do trabalho e na identificação de impedimentos, evitando sobrecarga de reuniões.

As revisões de sprint são direcionadas à validação das funcionalidades entregues e à coleta de feedback dos stakeholders, enquanto as retrospectivas são utilizadas para refletir sobre o processo de trabalho e identificar oportunidades de melhoria contínua. Essas adaptações garantem que os eventos do Scrum cumpram seu propósito, sem se tornarem burocráticos.

\section{Adaptação dos Artefatos}

Os artefatos do Scrum foram mantidos como instrumentos centrais de transparência e comunicação. O Product Backlog é utilizado como principal mecanismo de priorização e adaptação do escopo, sendo continuamente refinado com base no aprendizado obtido ao longo do projeto.

O Sprint Backlog é utilizado para dar visibilidade ao trabalho selecionado para cada sprint, permitindo o acompanhamento do progresso e a identificação precoce de desvios. Além disso, são utilizados artefatos complementares, como quadros visuais e registros de decisões, para apoiar a comunicação e o controle do projeto.

\section{Integração com Práticas Complementares}

Como parte do tailoring, foram incorporadas práticas complementares ao Scrum, quando consideradas necessárias para atender às demandas do projeto. Entre essas práticas destacam-se o uso de métricas simples de acompanhamento, a documentação mínima necessária para garantir rastreabilidade e a definição de critérios de aceitação claros para cada funcionalidade.

Essa integração busca equilibrar a flexibilidade dos métodos ágeis com a necessidade de controle e previsibilidade exigida pelo contexto institucional, reforçando o alinhamento com os domínios de desempenho do PMBOK 7.

\section{Limites e Restrições do Tailoring}

O tailoring adotado reconhece limitações inerentes ao projeto, como restrições de tempo, recursos humanos e maturidade organizacional. Nem todas as práticas ágeis são aplicadas em sua forma ideal, e algumas adaptações são necessárias para garantir viabilidade e alinhamento com o contexto acadêmico e institucional.

Essas restrições são tratadas de forma explícita, permitindo que o projeto mantenha coerência metodológica e foco em valor, mesmo diante de limitações práticas. Dessa forma, o tailoring adotado reforça o caráter pragmático da abordagem de gestão de projetos proposta neste trabalho.
